\clearpage\pagestyle{otMain}\setlength{\otBlocEspacement}{2.5mm}\fontsize{11}{15}\selectfont\corrigesChapitre{10}{Moyenne, médiane et étendue}
\begin{blocLivre}{cor-10-1}\begin{corrige}{1}{Un effectif total}\demarche{Additionner les effectifs avant de calculer la fréquence.}
Le total vaut $3+5+2=\res{10}$. La fréquence de 10 est $\frac5{10}=\res{0{,}5}$, soit 50\,\%.
\end{corrige}\end{blocLivre}
\begin{blocLivre}{cor-10-2}\begin{corrige}{2}{Une moyenne simple}\demarche{Additionner toutes les valeurs et compter les données.}
La somme vaut $6+9+10+15=40$. On divise par quatre valeurs : $40:4=\res{10}$.
\end{corrige}\end{blocLivre}
\begin{blocLivre}{cor-10-3}\begin{corrige}{3}{Une médiane impaire}\demarche{Ranger avant de chercher le rang central.}
La série rangée est 7, 8, 9, 12, 15. L’effectif est 5 : la 3e valeur est centrale. La médiane vaut $\res9$ : deux valeurs sont plus petites et deux plus grandes.
\end{corrige}\end{blocLivre}
\begin{blocLivre}{cor-10-4}\begin{corrige}{4}{Deux étendues}\demarche{Repérer le maximum et le minimum.}
Première série : $15-7=\res8$. Seconde série : $4-(-6)=\res{10}$. On soustrait toujours la plus petite valeur à la plus grande.
\end{corrige}\end{blocLivre}
\begin{blocLivre}{cor-10-5}\begin{corrige}{5}{Un effectif manquant}\demarche{Utiliser l’effectif total.}
L’effectif manquant vaut $25-(4+9+5)=\res7$. La fréquence de la note 10 est $\frac7{25}=\res{0{,}28}$, soit 28\,\%.
\end{corrige}\end{blocLivre}
\begin{blocLivre}{cor-10-6}\begin{corrige}{6}{Une moyenne pondérée}\demarche{Tenir compte des répétitions.}
L’effectif total vaut 10. La somme pondérée est $1\times5+4\times2+6\times3=31$. La moyenne est $\frac{31}{10}=\res{3{,}1}$. Moyenner seulement 1, 4 et 6 donnerait $3{,}67$ environ : un résultat faux.
\end{corrige}\end{blocLivre}
\begin{blocLivre}{cor-10-7}\begin{corrige}{7}{Lire un diagramme en bâtons}\demarche{Retrouver les rangs à partir des effectifs.}
Les effectifs sont 2, 4, 3, 2 et 1 pour 0, 1, 2, 3 et 4 livres. Moyenne : $\frac{0+4+6+6+4}{12}=\frac{20}{12}\approx\res{1{,}67}$ livre. Les 6e et 7e valeurs sont 1 et 2 : toute valeur entre 1 et 2 est une médiane, par exemple $\res{1{,}5}$. Étendue : $4-0=\res4$ livres.
\label{fig-cor-10-7}\figureGeo[corrige]{stat-livres}
\end{corrige}\end{blocLivre}
\begin{blocLivre}{cor-10-8}\begin{corrige}{8}{Une médiane paire}\demarche{Repérer les deux rangs centraux.}
Il y a six valeurs, déjà rangées. Les valeurs centrales sont aux rangs 3 et 4 : 8 et 10. Toute valeur comprise entre 8 et 10 est une médiane ; on choisit leur moyenne, $\res9$. Trois valeurs sont inférieures à 9 et trois supérieures.
\end{corrige}\end{blocLivre}
\begin{blocLivre}{cor-10-9}\begin{corrige}{9}{Un diagramme circulaire}\demarche{Passer de l’angle à l’effectif.}
Les valeurs vont de 0 à 5 : l’étendue vaut $\res5$. Un angle de $90^\circ$ représente $\frac{90}{360}=\frac14$ des élèves, soit 10 élèves. Les effectifs sont $\res{10,\ 16,\ 8,\ 4,\ 2}$ pour 0, 1, 2, 3 et 5. Moyenne : $\frac{0+16+16+12+10}{40}=\frac{54}{40}=\res{1{,}35}$.
\label{fig-cor-10-9}\figureGeo[corrige]{stat-fratrie}
\end{corrige}\end{blocLivre}
\begin{blocLivre}{cor-10-10}\begin{corrige}{10}{Des formules de tableur}\demarche{Associer chaque indicateur à sa fonction.}
\texttt{=MOYENNE(A2:A11)} affiche $\res{15{,}5}$ ; \texttt{=MEDIANE(A2:A11)} affiche $\res{15}$ ; \texttt{=MAX(A2:A11)-MIN(A2:A11)} affiche $30-9=\res{21}$ ; \texttt{=NB(A2:A11)} affiche $\res{10}$. Les valeurs rangées ont pour 5e et 6e valeurs 15 et 15.
\end{corrige}\end{blocLivre}
\begin{blocLivre}{cor-10-11}\begin{corrige}{11}{Ajouter une valeur}\demarche{Retrouver la somme à partir de la moyenne.}
La somme initiale vaut $5\times12=60$. La nouvelle somme vaut $60+18=78$, pour 6 élèves. La moyenne est $78:6=\res{13}$. Moyenner 12 et 18 donnerait un mauvais poids aux cinq premiers élèves.
\end{corrige}\end{blocLivre}
\begin{blocLivre}{cor-10-12}\begin{corrige}{12}{Une valeur manquante}\demarche{Utiliser la moyenne pour retrouver le total.}
La somme doit valoir $4\times12=48$. Donc $x=48-(8+11+14)=\res{15}$. La série rangée est 8, 11, 14, 15 : toute valeur entre 11 et 14 est une médiane, par exemple $\res{12{,}5}$.
\end{corrige}\end{blocLivre}
\begin{blocLivre}{cor-10-13}\begin{corrige}{13}{Une valeur extrême}\demarche{Calculer séparément les trois indicateurs.}
Pour A : moyenne 6, médiane 6, étendue 4. Pour B, la somme vaut 60 : moyenne $\res{12}$, médiane toujours $\res6$, étendue $\res{34}$. Modifier la seule plus grande valeur change fortement la moyenne et l’étendue, pas la médiane.
\end{corrige}\end{blocLivre}
\begin{blocLivre}{cor-10-14}\begin{corrige}{14}{Comparer deux classes}\demarche{Utiliser le sens de chaque indicateur.}
La classe A a une moyenne de $\frac{30}{10}=3$ livres : $\res{B}$ a la plus grande moyenne. Dans A, trois élèves empruntent au moins 4 livres. Dans B, la médiane 4 montre qu’au moins 6 élèves sur 12 empruntent au moins 4 livres : $\res{B}$ en compte le plus. L’étendue de A vaut 8, celle de B 6 : les valeurs de $\res{A}$ sont les plus dispersées.
\end{corrige}\end{blocLivre}
\begin{blocLivre}{cor-10-15}\begin{corrige}{15}{Interpréter une médiane}\demarche{Distinguer valeur centrale et effectif d’une valeur.}
Non : une médiane parle de position dans la série ordonnée, pas du nombre de valeurs égales à 15. Par exemple, dix durées de 10 min et dix de 20 min : toute valeur entre 10 et 20 est une médiane, donc 15 aussi, alors $\res{\text{aucune durée ne vaut 15 min}}$.
\end{corrige}\end{blocLivre}
